\documentclass{article}
\usepackage[UTF8]{ctex}
\usepackage{amsmath}
\usepackage{amssymb}
\usepackage{graphicx}
\usepackage{geometry}
\geometry{a4paper, margin=2cm}

\newcommand{\problemname}[1]{\begin{center}\Large\bfseries #1\end{center}}
\usepackage{enumitem}

\begin{document}
\problemname{实现蜗牛爬树问题}

\subsection*{题目描述}
一只蜗牛从 n 米高的大树底部向树顶攀爬。白天它向上爬 a 米，晚上会向下滑落 b 米；第二天白天再爬 a 米，晚上再滑落 b 米，如此往复。请计算按照这样的规律，蜗牛第几天可以爬到树顶。

\subsection*{输入格式}
输入只有一行，包含三个用空格分开的整数 n、a、b，分别表示树的高度、白天爬升的高度和晚上滑落的高度。保证 a > b，并且 n - a 是 (a - b) 的整数倍。所有数值都在 int 类型变量可以存储的范围内。

\subsection*{输出格式}
输出一个整数，表示蜗牛爬到树顶所需的天数。

\subsection*{样例输入}
\begin{verbatim}
12 3 2
\end{verbatim}

\subsection*{样例输出}
\begin{verbatim}
10
\end{verbatim}


\end{document}
